\begin{proposition}[5.9.2]
\label{IV.2.s5.9.2-fr}
\begin{enumerate}
\item[\textup{(i)}] El funtor
$\mathcal F\leadsto\mathcal H^0_{X/Z}(\mathcal F)$ es exacto por la izquierda.
\item[\textup{(ii)}] Si $\mathcal F$ es cuasicoherente, lo mismo ocurre con
$\mathcal H^0_{X/Z}(\mathcal F)$.
\end{enumerate}
\end{proposition}

La afirmación (i) resulta de la definición (5.9.1.1), del hecho de que
$(i_\alpha)_*$ es un funtor exacto por la izquierda y de que el límite
inductivo preserva la exactitud en la categoría de los $\mathcal O_X$-módulos.
La afirmación (ii) resulta de (\textbf I, 9.2.2) y del hecho de que un límite inductivo
de $\mathcal O_X$-módulos cuasicoherentes es cuasicoherente (\textbf I, 1.3.9).

\begin{remark}[5.9.3]
\label{IV.2.s5.9.3-fr}
Si $\mathcal F$ es una $\mathcal O_X$-\emph{álgebra}, lo mismo ocurre con
$\mathcal H^0_{X/Z}(\mathcal F)$ $(\mathbf{0}_{\mathrm I},4.2.4)$; en particular,
$\mathcal H^0_{X/Z}(\mathcal O_X)$ es una $\mathcal O_X$-álgebra
cuasicoherente y, para todo $\mathcal O_X$-módulo $\mathcal F$,
$\mathcal H^0_{X/Z}(\mathcal F)$ es un
$\mathcal H^0_{X/Z}(\mathcal O_X)$-módulo, que es cuasicoherente si
$\mathcal F$ es cuasicoherente (\textbf I, 9.6.1). Más concretamente,
supongamos que $X=\operatorname{Spec}(A)$, donde $A$ es \emph{íntegro} y noetheriano;
entonces $\mathcal H^0_{X/Z}(\mathcal O_X)$ es la $\mathcal O_X$-álgebra
$\widetilde B$, donde
\begin{equation}
B=\bigcap_{\mathfrak p\in X-Z}A_{\mathfrak p}.
\tag{5.9.3.1}\label{IV.2.s5.9.3.1-fr}
\end{equation}
Esto resulta, en efecto, de (5.9.1.2) y de (\textbf I, 8.2.1.1).
\end{remark}

\begin{proposition}[5.9.4]
\label{IV.2.s5.9.4-fr}
Sean $X$ un preesquema localmente noetheriano, $Z$ una parte de $X$ estable
por especialización, $X'$ un preesquema localmente noetheriano y
$f:X'\to X$ un morfismo plano. Entonces $Z'=f^{-1}(Z)$ es estable por
especialización y, para todo $\mathcal O_X$-módulo cuasicoherente $\mathcal F$,
se tiene un isomorfismo canónico
\begin{equation}
f^*(\mathcal H^0_{X/Z}(\mathcal F))
\xrightarrow{\sim}\mathcal H^0_{X'/Z'}(f^*(\mathcal F))
\tag{5.9.4.1}\label{IV.2.s5.9.4.1-fr}
\end{equation}
\end{proposition}

En efecto, con las notaciones de (5.9.1), $Z'_\alpha=f^{-1}(Z_\alpha)$ es
cerrado en $X'$ y $Z'$ es la unión de los $Z'_\alpha$; además,
$(i_\alpha)_{(X')}$ es la inyección canónica
$i'_\alpha:U'_\alpha\to X'$, si
$U'_\alpha=X'\setmin Z'_\alpha=f^{-1}(U_\alpha)$. Como $f$ es plano, se sabe
(2.3.1) que el homomorfismo canónico
$f^*((i_\alpha)_*(\mathcal F|U_\alpha))
\longrightarrow(i'_\alpha)_*(f^*(\mathcal F)|U'_\alpha)$
es biyectivo; como, para $\alpha\leq\beta$, el diagrama
\[
\xymatrix{
f^*((i_\alpha)_*(\mathcal F|U_\alpha))
  \ar[r]^{\sim}\ar[d] &
(i'_\alpha)_*(f^*(\mathcal F)|U'_\alpha)\ar[d]\\
f^*((i_\beta)_*(\mathcal F|U_\beta))
  \ar[r]^{\sim} &
(i'_\beta)_*(f^*(\mathcal F)|U'_\beta)
}
\]
\oldpage[IV-2]{111}
es conmutativo, al pasar al límite se obtiene, por tanto, un isomorfismo canónico
$\varinjlim_\alpha
(f^*((i_\alpha)_*(\mathcal F|U_\alpha)))
\xrightarrow{\sim}\mathcal H^0_{X'/Z'}(f^*(\mathcal F))$.

Pero, como el funtor $f^*$ conmuta con los límites inductivos
$(\mathbf{0}_{\mathrm I},4.3.2)$, esto da por definición el isomorfismo
(5.9.4.1) buscado.
